My weld-seam HMI taught me to extract a structured signal from images, while CR5 deployment showed that a perceptual input becomes useful only when sequential decisions and controller timing remain coherent. I now need to connect these layers through deeper training in embodied perception, sequential decision making, and control validation. If the current curriculum continues into my intake, I would choose the Autonomous Driving stream within CityUHK's MSc in Artificial Intelligence because it frames that perception-decision-control problem around an embodied platform.

With those offerings available, CS5493 Topics in Autonomous Driving would provide the application context for examining the whole chain, while SDSC6007 Dynamic Programming and Reinforcement Learning would deepen how I formulate decisions over time. As a Group II elective, CS5187 Vision and Image could strengthen the visual representations feeding those decisions. If I chose the optional, mutually exclusive project alternative and received approval, the six-credit CS6522 Project in Autonomous Driving could let me test how visual uncertainty and policy timing propagate to control behavior. I would bring experience diagnosing boundaries between perception software and robot execution, then develop toward embodied-AI R\&D in which sensing, decisions, and physical response are validated together.
